\maketableofcontents

\chapter*{ВВЕДЕНИЕ}
\addcontentsline{toc}{chapter}{ВВЕДЕНИЕ}

В современном мире сервисы краткосрочной аренды автомобилей становятся одним из распространенных способов получения транспорта без необходимости владения личным автомобилем. Пользователю важно быстро выбрать подходящую машину, указать даты аренды, оформить бронирование и получить подтверждение оплаты через единый веб-интерфейс. Для администратора такой системы важны управление учетными записями, контроль событий и получение статистики по действиям пользователей.

Разработка подобной системы требует не только реализации пользовательского интерфейса, но и построения распределенной архитектуры. Отдельные компоненты должны отвечать за авторизацию, учет автомобилей, бронирование, платежи, сбор статистики и агрегацию запросов. Такой подход позволяет изолировать ответственность сервисов, упростить сопровождение и обеспечить независимое развертывание компонентов.

Особое значение в распределенной системе имеет безопасность взаимодействия. Все запросы между сервисами должны выполняться с использованием токенов, выданных Identity Provider по протоколу OpenID Connect. Для проверки подлинности токенов используется JWT и набор ключей JWKs, что позволяет сервисам валидировать запросы без обращения к сервису авторизации при каждом пользовательском действии.

Целью курсовой работы является разработка прототипа системы бронирования автомобилей на выбранные даты на базе веб-интерфейса.

Для достижения поставленной цели необходимо решить следующие задачи:

\begin{itemize}
  \item описать предметную область и функциональные требования к системе бронирования автомобилей;
  \item спроектировать микросервисную архитектуру системы;
  \item реализовать сервис пользовательского интерфейса;
  \item реализовать собственный сервис авторизации и данных пользовательских аккаунтов с поддержкой OpenID Connect;
  \item реализовать сервисы автомобилей, бронирования, платежей, статистики и агрегирования запросов;
  \item организовать асинхронную передачу событий через Apache Kafka;
  \item обеспечить авторизацию межсервисных запросов на основе JWT-токенов;
  \item подготовить развертывание системы в Kubernetes с использованием Helm charts и Ingress;
  \item настроить сборку и развертывание сервисов через CI/CD.
\end{itemize}
